\chapter*{Chapter 0: The Machine Underneath}
\addcontentsline{toc}{chapter}{Chapter 0: The Machine Underneath}

There is a proof underneath this book, and it is not made of prose. Every claim the body makes
in its metaphysical voice restates a step in a construction written in a proof assistant and
checked, line by line, by a computer. The construction is real, it is finite, and it compiles.
The body never mentions it. This chapter, alone with the preface, does.

What the body gives is an armwave. That is not an insult to it; it is the exact description of
what it is. The body describes a process --- the mark, the receipt, the residue, the boundary,
the bit --- faithfully and in order, but it describes the process, not the mechanism that
performs it. An armwave, done well, is a true account of what a machine does without being the
machine. The body is the true account. The machine is underneath.

This chapter describes the mechanism, and in describing the mechanism it does the one thing the
armwave cannot: it describes how the proof proves itself. The mechanism is not a fixed procedure
the author ran once and wrote down. It is a construction that, run, generates its own proof. To
say how it works is to say how the book is established --- and the book's content is exactly that
proof. Describe the machine precisely and you have described how the book proves itself.

The mechanism turns on a single classical step. Everywhere else the construction is
constructive: it builds what it claims, and the building is the evidence. At one point, and one
point only, it reaches past the constructive, to the axiom of choice --- the assertion that a
choice can be made without a rule for making it. The body calls this the needle, the one decision
the apparatus sits in rather than derives. The machine treats it differently. It does not pick an
algorithm to make the choice. It inspects the algorithms that could.

A choice with no rule is a choice that many rules could imitate, and the machine considers them
as a field: the possible algorithms the axiom of choice might employ to do what it does without
saying how. It does not run them for their outputs. It encodes each by what it costs the proof
assistant to elaborate it --- the time the assistant spends bringing that algorithm to a checked
form. Elaboration time is the measurement the body has been taking all along, under the name of
the heartbeat; here it is the very code in which the candidate algorithms are written down. Each
possible way of making the unmade choice becomes a finite elaboration cost.

The proof assistant then builds a program out of these encodings. It staples them together ---
composes the elaboration-time codes into one assembled object, the way the construction stapled
its tower of instances into a single synthesized whole. The program is not written by hand and
handed to the assistant; it is assembled by the assistant, out of the encoded algorithms, as the
construction elaborates.

What gets built is, itself, a proof assistant. Not the whole of one --- it does not verify
arbitrary mathematics --- but a proof assistant insofar as the techniques the proof uses: it can
carry out, on its own, the checking the construction needs. The assistant has built an assistant,
restricted to the proof's own methods. This is the strange turn the chapter has been approaching,
and it is why the book leans on Hofstadter where it does: a formal system, to refer to itself
without standing outside itself, builds inside itself a smaller image of itself.

The last step closes the loop. The original proof assistant simulates a computer --- a finite
machine, modeled inside the assistant --- and runs the built program on it. The run is the proof.
Not a description of the proof, not a promise that one exists: the execution of the assembled
assistant, on the simulated machine, generates the proof the book restates. The assistant builds
an assistant, sets it running on a computer the assistant pretends into being, and the running is
the thing.

So the book is an armwave because the mechanism is this, and prose is not this. The body can
describe the process the machine performs --- it does, carefully, for twenty-one chapters --- but
the process described is the shadow the machine casts, not the loop the machine runs. The loop is
here, stated once: a proof assistant inspects the algorithms a free choice could employ, encodes
them by elaboration time, staples them into a program that is itself a proof assistant for the
proof's techniques, and runs that program on a computer it simulates --- and the running
generates the proof.

From here the book does not mention any of this again. The next page is the first mark, and from
the first mark to the last compass needle the book speaks only the theory --- the marks, the
residues, the coin on the table --- as though the machine underneath had never been built. It was
built. It compiles. The reader who wants to know what the armwave waves at has, in this chapter
and the preface, the whole of the admission; everything after is the theory, told straight.
